\section{Related Work}

Here, we review works in the field of GraphRAG and visual Analysis for LLM-based Applications, which are closely related to our study.

\subsection{GraphRAG}

GraphRAG is an emerging powerful retrieval-augmented generation paradigm that leverages knowledge graph as an intermediate representation to enable more precise and comprehensive retrieval for LLM generation, which has recently shown great potential in diverse application areas, such as medicine\cite{DALK}, E-commerce\cite{ecommerce}, intelligence analysis\cite{Ranade_2023}, and software engineering\cite{alhanahnah2024depsragagenticreasoningplanning}. While all GraphRAG approaches utilize graphs as a crucial intermediate knowledge representation, the underlying design space is vast and has attracted extensive research attention. For the graph construction phase, various indexing strategies (e.g. basic graph indexing\cite{gutiérrez2025hipporagneurobiologicallyinspiredlongterm}, text indexing\cite{li2024unioqaunifiedframeworkknowledge}, vector indexing\cite{he2024g}, and hybrid indexing\cite{sarmah2024hybridrag}) have been explored to facilitate downstream retrievals. For the graph retrieval phase, different types of recalls (e.g. nodes\cite{li2024graphneuralnetworkenhanced}, relationships\cite{li2024unioqaunifiedframeworkknowledge}, subgraphs\cite{edge2024local}, and paths\cite{ma2024think}) can be retrieved. Moreover, the retrieval paradigm can vary from efficient and cost-effective one-time retrieval \cite{he2024g, gutiérrez2025hipporagneurobiologicallyinspiredlongterm} to more sophisticated and adaptive iterative retrieval \cite{ma2024think, sun2023thinkongraph}.

While extensive work has focused on optimizing GraphRAG's process design, none has addressed how to assist developers in analyzing the underlying complex processes—a crucial aspect for practical GraphRAG applications. In this work, we propose a visual analytics framework that enables users to better understand and analyze GraphRAG processes, facilitating its effective deployment in real-world scenarios.

\subsection{Visual Analysis for LLM-based Applications}
Based on the interaction patterns of LLM invocations, existing visualization systems for analyzing LLM behavior can be categorized into two main directions: single-shot LLM invocation analysis and iterative LLM invocation analysis.

For single-shot LLM invocation analysis, researchers have developed various visualization systems to analyze individual LLM responses. PromptIDE \cite{strobelt2022interactive} enables users to experiment with prompt variations and visualize prompt performance for ad-hoc task adaptation. LLM Comparator \cite{kahng2024llm} facilitates side-by-side evaluation of LLM responses through interactive visual analytics, helping users understand performance differences between models. To analyze specific LLM capabilities, some work focused on specialized analysis tasks. For example, researchers proposed visualization techniques to analyze text style transfer \cite{BrathVisualizingLT}, revealing how LLMs apply different stylistic features in their responses. In terms of security analysis, JailbreakLens \cite{feng2024jailbreaklens} provides a visual analytics system for analyzing jailbreak attacks against LLMs, enabling users to evaluate model vulnerabilities and defensive capabilities.

For iterative LLM invocation analysis, existing work has explored how to visualize and analyze complex interaction patterns between multiple LLM calls. Sensecape \cite{suh2023sensecape} supports multilevel exploration and sensemaking with LLMs by enabling users to manage information complexity through hierarchical organization. ChainForge \cite{arawjo2024chainforge} provides a visual toolkit for prompt engineering and hypothesis testing across multiple LLM responses. Recent work has also focused on visually analyzing LLM-based agent systems. AgentLens \cite{lu2024agentlens} visualizes the dynamic evolution of agent behaviors in LLM-based autonomous systems, while AgentCoord \cite{pan2024agentcoordvisuallyexploringcoordination} helps users explore and design coordination strategies for LLM-based multi-agent collaboration. \cite{liu2023sprout} breaks down the programming tutorial creation process into actionable steps and visualizes the tree-of-thought exploration process.

Extending this research line, we aim to develop an interactive visual analysis framework for analyzing complex GraphRAG processes with diverse and intensive iterative LLM invocations.